% ============================================================
% CHƯƠNG 4: THỰC NGHIỆM VÀ ĐÁNH GIÁ
% ============================================================
\chapter{Thực nghiệm và đánh giá}

\section{Thiết lập thực nghiệm}

\subsection{Bộ dữ liệu đánh giá}
% TODO:
% - 50 bài báo từ 4 trang: VnExpress, Tuổi Trẻ, Dân Trí, Thanh Niên
% - 5 chủ đề: thời sự, pháp luật, kinh tế, giáo dục, văn hóa (10 bài/chủ đề)
% - Tiêu chí chọn bài: đa dạng độ dài, chủ đề, nguồn
% Bảng 4.1: Thống kê bộ dữ liệu đánh giá

\subsection{Các mô hình được đánh giá}
% TODO:
% - GPT-4o (OpenAI) - baseline đa ngôn ngữ
% - GPT-4o-mini (OpenAI) - phiên bản nhỏ hơn
% - ViT5-large (VietAI) - chuyên biệt tiếng Việt
% - Thảo luận: PhoGPT, Vistral chưa triển khai được (hosting constraints)
% Bảng 4.2: Thông tin các mô hình

\subsection{Các metrics đánh giá}
% TODO:
% - ROUGE-1, ROUGE-2, ROUGE-L
% - BLEU
% - BERTScore (PhoBERT-based)
% - Compression Rate
% - Latency (ms)

\subsection{Môi trường thực nghiệm}
% TODO:
% - Backend: Next.js trên localhost
% - BERTScore service: FastAPI trên HuggingFace Spaces
% - ViT5: HuggingFace Inference API
% - GPT-4o/4o-mini: OpenAI API

\section{Kết quả thực nghiệm}

\subsection{So sánh tổng quan các mô hình}
% TODO:
% Bảng 4.3: Kết quả trung bình các metrics theo model
% - ROUGE-1, ROUGE-2, ROUGE-L, BLEU, BERTScore, Compression, Latency
% - So sánh GPT-4o vs GPT-4o-mini vs ViT5
% Hình 4.1: Biểu đồ radar so sánh các model

\subsection{Phân tích theo chủ đề}
% TODO:
% Bảng 4.4: Kết quả theo chủ đề (thời sự)
% Bảng 4.5: Kết quả theo chủ đề (pháp luật)
% Bảng 4.6: Kết quả theo chủ đề (kinh tế)
% Bảng 4.7: Kết quả theo chủ đề (giáo dục)
% Bảng 4.8: Kết quả theo chủ đề (văn hóa)
% Hình 4.2: Biểu đồ so sánh BERTScore theo chủ đề

\subsection{Phân tích BERTScore chi tiết}
% TODO:
% - ViT5 đạt BERTScore cao hơn GPT-4o: phân tích nguyên nhân
% - PhoBERT embeddings ưu tiên token overlap tiếng Việt
% - ViT5 sinh tóm tắt extractive hơn → BERTScore cao hơn
% - GPT-4o sinh tóm tắt abstractive hơn → đọc tự nhiên hơn nhưng BERTScore thấp hơn
% Hình 4.3: Phân bố BERTScore theo model

\subsection{Phân tích Compression Rate và Latency}
% TODO:
% - GPT-4o: tóm tắt dài hơn, latency ổn định
% - ViT5: tóm tắt ngắn hơn, latency phụ thuộc HF Inference API
% Hình 4.4: Scatter plot Compression Rate vs BERTScore
% Hình 4.5: Box plot Latency theo model

\section{Đánh giá cơ chế Routing}

\subsection{Hiệu quả phân loại độ phức tạp}
% TODO:
% - Accuracy của phân loại simple/medium/complex
% - Ví dụ minh họa routing decisions

\subsection{So sánh routing tự động vs manual}
% TODO:
% - Routing tự động có chọn đúng model tối ưu không?
% - Trade-off giữa chất lượng và chi phí API

\section{Đánh giá module kiểm chứng sự kiện}

\subsection{Chất lượng kiểm chứng}
% TODO:
% - Đánh giá định tính: mẫu fact-check results
% - Độ chính xác của trust scoring
% - Chất lượng nguồn tìm kiếm từ Tavily

\subsection{Hạn chế}
% TODO:
% - Phụ thuộc chất lượng tìm kiếm
% - Không thể kiểm chứng tin tức quá mới (chưa có nguồn)
% - Bias của LLM trong đánh giá

\section{Thảo luận về thách thức triển khai mô hình tiếng Việt}

% TODO:
% - PhoGPT-4B: yêu cầu GPU (T4 16GB), ZeroGPU cần HF Pro
% - Vistral (Viet-Mistral): tương tự, GPU requirement
% - So sánh chi phí hosting: HF Spaces vs Modal.com vs self-hosted
% - Hàm ý: model chuyên biệt tiếng Việt cần hạ tầng GPU, rào cản cho nghiên cứu

\section{Phân tích key findings}

% TODO:
% Finding 1: ViT5 > GPT-4o về BERTScore nhưng < về tính tự nhiên
% Finding 2: Model chuyên biệt (fine-tuned) cho kết quả ngữ nghĩa tốt hơn
%            trên tác vụ cụ thể, dù kích thước nhỏ hơn nhiều
% Finding 3: Cơ chế routing giảm chi phí API mà vẫn đảm bảo chất lượng
% Finding 4: BERTScore với PhoBERT phản ánh chất lượng ngữ nghĩa tiếng Việt
%            tốt hơn ROUGE/BLEU
% Finding 5: Search-augmented fact-checking khả thi cho tin tức tiếng Việt

\section{Tổng kết chương}
% TODO: Tóm tắt kết quả chính và nhận xét
